% GENERATED FILE. Edit canonical lessons, then run npm run generate:books.
% canonical-source-hash: fnv1a64:87e7d4724a1ed012
% canonical-lessons: JA-C269-dorobou, JA-C269-nusumu, JA-C269-yuube, JA-C269-hoeru, JA-C269-sawagu

\chapter{A Thief in the Night}
\label{ch:ja-a-thief-in-the-night}
\hlchaptermodality{269}

\begin{chapteropening}
\textbf{By the end of this chapter:} \emph{I can say what a thief did and what we heard.}

Complete the last lesson of chapter 269: I can say what a thief did and what we heard.
\end{chapteropening}

\section[\texorpdfstring{\ja{どろぼう} (dorobō)}{どろぼう (dorobō)}]{\ja{どろぼう} (dorobō) — a thief; a burglar}

\label{lesson:JA-C269-dorobou}



\begin{quote}
Before the new one: say the Japanese for to help each other, then the Japanese for safe.
\end{quote}

\subsection*{You'll want to know: \ja{どろぼう}}
\textbf{\ja{どろぼう}} — \emph{dorobō} — \textquotedblleft{}a thief; a burglar\textquotedblright{}.

\begin{quote}
\textbf{\ja{どろぼうが} \ja{うちに} \ja{はいりました。}} — \emph{dorobō ga uchi ni hairimashita} — \textquotedblleft{}A burglar broke into our house.\textquotedblright{}
\end{quote}

\begin{cousinweb}
\textbf{\ja{どろぼう}} (\emph{dorobō}) is a mixed word, native and Sino-Japanese: \emph{doro}, \textquotedblleft{}mud\textquotedblright{}, is native Japanese, and \emph{bō}, \textquotedblleft{}stick\textquotedblright{}, is a Chinese character read with its on'yomi, a reading borrowed from Chinese (Middle Chinese \emph{baewngX}) (JMdict, KANJIDIC2; Middle Chinese in Baxter's transcription). How the two came together is not settled.
\end{cousinweb}

\begin{grammarlens}[title={What you've built}]
One more word for telling what happened.
\end{grammarlens}

\subsection*{Your turn}
\begin{itemize}
\raggedright
  \item \emph{Say it:} \emph{dorobō}
  \item \emph{Say it:} \emph{dorobō}, once more
  \item \emph{Say it:} \emph{dorobō ga uchi ni hairimashita} — A burglar broke into our house.
  \item \emph{Recall:} say the Japanese for to help each other, then the Japanese for safe, then say \emph{dorobō} again
\end{itemize}

\begin{tcolorbox}[breakable,colback=teal!4,colframe=teal!35!black,title={Before you move on}]
What does \ja{どろぼう} mean? (\textquotedblleft{}A thief; a burglar\textquotedblright{}.) Say it once more.
\end{tcolorbox}
\section[\texorpdfstring{\ja{ぬすむ} (nusumu)}{ぬすむ (nusumu)}]{\ja{ぬすむ} (nusumu) — to steal}

\label{lesson:JA-C269-nusumu}



\begin{quote}
Before the new one: say the Japanese for safe, then the Japanese for a thief.
\end{quote}

\subsection*{You'll want to know: \ja{ぬすむ}}
\textbf{\ja{ぬすむ}} — \emph{nusumu} — \textquotedblleft{}to steal\textquotedblright{}.

\begin{quote}
\textbf{\ja{どろぼうは} \ja{なにも} \ja{ぬすみませんでした。}} — \emph{dorobō wa nanimo nusumimasen deshita} — \textquotedblleft{}The burglar took nothing.\textquotedblright{}
\end{quote}

Polite \textbf{\ja{ぬすみます}}, \emph{nusumimasu}; a godan verb. Past (polite): \textbf{\ja{どろぼうが} \ja{さいふを} \ja{ぬすみました。}}, \emph{dorobō ga saifu o nusumimashita}.

\begin{cousinweb}
\textbf{\ja{ぬすむ}} (\emph{nusumu}) is a native Japanese word (wago), inherited rather than borrowed from Chinese; its kanji was attached later as a kun'yomi (JMdict).
\end{cousinweb}

\begin{grammarlens}[title={What you've built}]
One more word for telling what happened.
\end{grammarlens}

\subsection*{Your turn}
\begin{itemize}
\raggedright
  \item \emph{Say it:} \emph{nusumu}
  \item \emph{Say it:} \emph{nusumu}, once more
  \item \emph{Say it:} \emph{dorobō wa nanimo nusumimasen deshita} — The burglar took nothing.
  \item \emph{Recall:} say the Japanese for safe, then the Japanese for a thief, then say \emph{nusumu} again
\end{itemize}

\begin{tcolorbox}[breakable,colback=teal!4,colframe=teal!35!black,title={Before you move on}]
What does \ja{ぬすむ} mean? (\textquotedblleft{}To steal\textquotedblright{}.) Say it once more.
\end{tcolorbox}
\section[\texorpdfstring{\ja{ゆうべ} (yūbe)}{ゆうべ (yūbe)}]{\ja{ゆうべ} (yūbe) — last night}

\label{lesson:JA-C269-yuube}



\begin{quote}
Before the new one: say the Japanese for a thief, then the Japanese for to steal.
\end{quote}

\subsection*{You'll want to know: \ja{ゆうべ}}
\textbf{\ja{ゆうべ}} — \emph{yūbe} — \textquotedblleft{}last night\textquotedblright{}.

\begin{quote}
\textbf{\ja{ゆうべ} \ja{けいさつから} \ja{でんわが} \ja{ありました。}} — \emph{yūbe keisatsu kara denwa ga arimashita} — \textquotedblleft{}Last night the police phoned us.\textquotedblright{}
\end{quote}

\begin{cousinweb}
\textbf{\ja{ゆうべ}} (\emph{yūbe}) is a native Japanese word (wago), inherited rather than borrowed from Chinese; its kanji was attached later as a kun'yomi (JMdict).
\end{cousinweb}

\begin{grammarlens}[title={What you've built}]
One more word for telling what happened.
\end{grammarlens}

\subsection*{Your turn}
\begin{itemize}
\raggedright
  \item \emph{Say it:} \emph{yūbe}
  \item \emph{Say it:} \emph{yūbe}, once more
  \item \emph{Say it:} \emph{yūbe keisatsu kara denwa ga arimashita} — Last night the police phoned us.
  \item \emph{Recall:} say the Japanese for a thief, then the Japanese for to steal, then say \emph{yūbe} again
\end{itemize}

\begin{tcolorbox}[breakable,colback=teal!4,colframe=teal!35!black,title={Before you move on}]
What does \ja{ゆうべ} mean? (\textquotedblleft{}Last night\textquotedblright{}.) Say it once more.
\end{tcolorbox}
\section[\texorpdfstring{\ja{ほえる} (hoeru)}{ほえる (hoeru)}]{\ja{ほえる} (hoeru) — to bark; to howl}

\label{lesson:JA-C269-hoeru}



\begin{quote}
Before the new one: say the Japanese for to steal, then the Japanese for last night.
\end{quote}

\subsection*{You'll want to know: \ja{ほえる}}
\textbf{\ja{ほえる}} — \emph{hoeru} — \textquotedblleft{}to bark; to howl\textquotedblright{}.

\begin{quote}
\textbf{\ja{もし} \ja{どろぼうが} \ja{きたら、いぬが} \ja{ほえるでしょう。}} — \emph{moshi dorobō ga kitara, inu ga hoeru deshō} — \textquotedblleft{}If a burglar came, the dog would bark.\textquotedblright{}
\end{quote}

Polite \textbf{\ja{ほえます}}, \emph{hoemasu}; an ichidan verb. Past (polite): \textbf{\ja{いぬが} \ja{ほえました。}}, \emph{inu ga hoemashita}.

\begin{cousinweb}
\textbf{\ja{ほえる}} (\emph{hoeru}) is a native Japanese word (wago), inherited rather than borrowed from Chinese; its kanji was attached later as a kun'yomi (JMdict).
\end{cousinweb}

\begin{grammarlens}[title={What you've built}]
One more word for telling what happened.
\end{grammarlens}

\subsection*{Your turn}
\begin{itemize}
\raggedright
  \item \emph{Say it:} \emph{hoeru}
  \item \emph{Say it:} \emph{hoeru}, once more
  \item \emph{Say it:} \emph{moshi dorobō ga kitara, inu ga hoeru deshō} — If a burglar came, the dog would bark.
  \item \emph{Recall:} say the Japanese for to steal, then the Japanese for last night, then say \emph{hoeru} again
\end{itemize}

\begin{tcolorbox}[breakable,colback=teal!4,colframe=teal!35!black,title={Before you move on}]
What does \ja{ほえる} mean? (\textquotedblleft{}To bark; to howl\textquotedblright{}.) Say it once more.
\end{tcolorbox}
\section[\texorpdfstring{\ja{さわぐ} (sawagu)}{さわぐ (sawagu)}]{\ja{さわぐ} (sawagu) — to make a noise; to make a fuss}

\label{lesson:JA-C269-sawagu}



\begin{quote}
Before the new one: say the Japanese for last night, then the Japanese for to bark.
\end{quote}

\subsection*{You'll want to know: \ja{さわぐ}}
\textbf{\ja{さわぐ}} — \emph{sawagu} — \textquotedblleft{}to make a noise; to make a fuss\textquotedblright{}.

\begin{quote}
\textbf{\ja{となりの} \ja{へやで} \ja{おとこが} \ja{さわいでいました。}} — \emph{tonari no heya de otoko ga sawaide imashita} — \textquotedblleft{}A man was making a racket in the next room.\textquotedblright{}
\end{quote}

Polite \textbf{\ja{さわぎます}}, \emph{sawagimasu}; a godan verb. Past (polite): \textbf{\ja{こどもが} \ja{さわぎました。}}, \emph{kodomo ga sawagimashita}.

\begin{cousinweb}
\textbf{\ja{さわぐ}} (\emph{sawagu}) is a native Japanese word (wago), inherited rather than borrowed from Chinese; its kanji was attached later as a kun'yomi (JMdict).
\end{cousinweb}

\begin{grammarlens}[title={What you've built}]
One more word for telling what happened.
\end{grammarlens}

\subsection*{Your turn}
\begin{itemize}
\raggedright
  \item \emph{Say it:} \emph{sawagu}
  \item \emph{Say it:} \emph{sawagu}, once more
  \item \emph{Say it:} \emph{tonari no heya de otoko ga sawaide imashita} — A man was making a racket in the next room.
  \item \emph{Recall:} say the Japanese for last night, then the Japanese for to bark, then say \emph{sawagu} again
\end{itemize}

\begin{tcolorbox}[breakable,colback=teal!4,colframe=teal!35!black,title={Before you move on}]
What does \ja{さわぐ} mean? (\textquotedblleft{}To make a noise; to make a fuss\textquotedblright{}.) Say it once more.
\end{tcolorbox}
